\chapter{Aplicaciones de la integral}

La integral definida, más allá de su interpretación como área bajo una curva,
es una herramienta fundamental para modelar y resolver problemas geométricos,
físicos y de ingeniería. En este capítulo desarrollaremos las aplicaciones
clásicas: cálculo de áreas de regiones planas, volúmenes de sólidos de
revolución, longitud de arco, áreas de superficies de revolución y
aplicaciones físicas como trabajo, centro de masa y momento de inercia. El
hilo conductor de todo el capítulo es el método de
\textbf{rebanar, aproximar e integrar}: identificar la cantidad diferencial
$dQ$ apropiada, sumar infinitesimalmente las contribuciones y reconocer la
suma límite como una integral definida.

En general, un protocolo práctico puede resumirse así para cada aplicación:

\begin{rem}
\begin{description}
  \item[Áreas.] Grafique primero. Si las curvas se expresan mejor como
    $x=g(y)$, integre respecto a $y$. Cuando haya varios cruces, sume
    integrales en subintervalos.

  \item[Volúmenes.]
    \begin{itemize}
      \item \emph{Discos/arandelas:} el rectángulo representativo es
        \textbf{perpendicular} al eje de rotación.
      \item \emph{Cascarones:} el rectángulo representativo es
        \textbf{paralelo} al eje de rotación.
      \item Para ejes desplazados, ajuste el radio sumando o restando el
        desplazamiento.
    \end{itemize}

  \item[Longitud de arco.] Verifique que $f'$ sea continua; si la integral
    resultante no es elemental, use integración numérica (Simpson, Trapecios
    o Riemann).

  \item[Superficies.] Recuerde el factor $2\pi\times\text{radio}\times ds$;
    el radio es la distancia del punto al eje de rotación.

  \item[Trabajo.] Identifique $F(x)$ e integre. Para bombeo, exprese la
    distancia de elevación en función de la variable de integración.

  \item[Centro de masa.] Para regiones entre dos curvas:
    $M_y=\rho\int_a^b x[f(x)-g(x)]\,dx$ y
    $M_x=\frac{\rho}{2}\int_a^b([f(x)]^2-[g(x)]^2)\,dx$.

  \item[Fuerza hidrostática.] Use $F=\int\rho g h\,w(h)\,dh$; exprese la
    anchura $w(h)$ en función de la profundidad $h$.
\end{description}
\end{rem}

\begin{theorem}[Principio unificador de las aplicaciones de la integral]
Cualquier cantidad acumulativa $Q$ que satisfaga:
\begin{enumerate}
  \item Puede aproximarse en subintervalos: $Q\approx\sum g(x_i^*)\Delta x$.
  \item La aproximación converge cuando $\Delta x\to 0$.
\end{enumerate}
puede calcularse exactamente como $Q=\int_a^b g(x)\,dx$. Este principio
unifica todas las aplicaciones presentadas en el capítulo: el área, el
volumen, la longitud de arco, el trabajo, los momentos y la fuerza
hidrostática son, en esencia, límites de sumas de Riemann con distintos
integrandos $g(x)$.
\end{theorem}

\begin{rem}
La potencia del cálculo integral radica en transformar problemas
geométricos y físicos complejos en operaciones algebraicas sistemáticas.
Dominar estas aplicaciones requiere práctica en tres pasos fundamentales:
\begin{enumerate}
\item Identificar el elemento diferencial $dQ$ adecuado para el problema,
\item Expresar $dQ$ en términos de una sola variable de integración, y
\item Determinar los límites de integración a partir de la geometría del
problema.
\end{enumerate}
\end{rem}

% ============================================================
\section{Cálculo de áreas de regiones planas}
% ============================================================

\begin{definition}[Área bajo una curva]
Sea $f:[a,b]\to\mathbb{R}$ una función continua y no negativa. El
\textbf{área} de la región acotada por la gráfica de $f$, el eje $x$ y las
rectas verticales $x=a$, $x=b$ se define como
\[
  A = \int_a^b f(x)\,dx.
\]
Si $f(x)\leq 0$ en $[a,b]$, el área geométrica es
$A = -\int_a^b f(x)\,dx = \int_a^b |f(x)|\,dx$.
\end{definition}

\begin{rem}
La integral definida calcula el \textbf{área algebraica} (con signo),
mientras que el área geométrica siempre es no negativa. Para obtener el
área real entre una curva y el eje $x$ cuando $f$ cambia de signo, debe
dividirse el intervalo en partes donde $f$ no cambie de signo e integrar
el valor absoluto en cada parte:
\[
  A = \int_a^b |f(x)|\,dx
      = \int_{\{f\ge 0\}} f(x)\,dx \;-\; \int_{\{f\le 0\}} f(x)\,dx.
\]
\end{rem}

\begin{example}[Área entre la curva y el eje $x$ con cambio de signo]
Calcule el área de la región acotada por $y = x^3 - x$ y el eje $x$ en
$[-1, 1]$.

\begin{myproof}
\textbf{Paso 1: Análisis gráfico y región.}
Las raíces de $x^3-x = x(x-1)(x+1)$ en $[-1,1]$ son $x=-1, 0, 1$.
La función es positiva en $(-1,0)$ (región azul) y negativa en $(0,1)$
(región roja). Para obtener el área geométrica debemos integrar el valor
absoluto en cada subintervalo.
\begin{center}
% ---------------------------------------------------------------
% CORRECCIÓN: se usan trayectorias nombradas + fillbetween para
%   evitar que los rellenos se superpongan o tapen entre sí.
%   Región positiva en azul, región negativa en rojo.
% ---------------------------------------------------------------
\begin{tikzpicture}[scale=0.9]
\begin{axis}[
  axis lines  = middle,
  xmin = -1.4,  xmax = 1.4,
  ymin = -0.52, ymax = 0.52,
  xlabel = {$x$}, ylabel = {$y$},
  xtick = {-1, -0.5, 0, 0.5, 1},
  ytick = {-0.4, -0.2, 0, 0.2, 0.4},
  grid  = major,
  grid style = {line width=0.3pt, draw=gray!25},
  tick style = {thick},
  tick label style = {font=\small},
  every axis/.append style = {font=\small}
]
%--- trayectorias invisibles para el relleno ---
\addplot[name path=fpos, domain=-1:0, samples=80, draw=none] {x^3 - x};
\addplot[name path=ejep, domain=-1:0, samples=2,  draw=none] {0};
\addplot[name path=fneg, domain= 0:1, samples=80, draw=none] {x^3 - x};
\addplot[name path=ejen, domain= 0:1, samples=2,  draw=none] {0};
%--- rellenos sin superposición ---
\addplot[fill=blue!25, opacity=0.80] fill between[of=fpos and ejep];
\addplot[fill=red!25,  opacity=0.80] fill between[of=ejen and fneg];
%--- curva visible encima del relleno ---
\addplot[domain=-1.3:1.3, samples=200, blue!80!black, thick] {x^3 - x};
%--- etiquetas de signo ---
\node[blue!70!black, font=\footnotesize\bfseries]
     at (axis cs:-0.5,  0.20) {$f>0$};
\node[red!70!black,  font=\footnotesize\bfseries]
     at (axis cs: 0.5, -0.20) {$f<0$};
\end{axis}
\end{tikzpicture}
\end{center}
\textbf{Paso 2: Elemento diferencial.} $dA = |f(x)|\,dx$.

\textbf{Paso 3: Límites de integración.} Dividimos en $[-1,0]$ y $[0,1]$.

\textbf{Paso 4: Evaluación.}
\begin{align*}
  A &= \int_{-1}^{0}(x^3-x)\,dx + \left|\int_{0}^{1}(x^3-x)\,dx\right| \\
    &= \left[\frac{x^4}{4}-\frac{x^2}{2}\right]_{-1}^{0}
       + \left|\left[\frac{x^4}{4}-\frac{x^2}{2}\right]_{0}^{1}\right| \\
    &= \left(0 - \Bigl(\tfrac{1}{4}-\tfrac{1}{2}\Bigr)\right)
       + \left|\tfrac{1}{4}-\tfrac{1}{2}\right| \\
    &= \tfrac{1}{4} + \tfrac{1}{4} = \boxed{\tfrac{1}{2}}.
\end{align*}
\end{myproof}
\end{example}

\begin{theorem}[Área entre dos curvas]
Sean $f$ y $g$ funciones continuas en $[a,b]$ tales que $f(x)\geq g(x)$
para todo $x\in[a,b]$. El área de la región acotada por ambas gráficas y
las rectas $x=a$, $x=b$ es
\[
  A = \int_a^b \bigl[f(x)-g(x)\bigr]\,dx.
\]
\end{theorem}

\begin{proof}
Dividimos $[a,b]$ en $n$ subintervalos de ancho $\Delta x=(b-a)/n$.
En cada subintervalo $[x_{i-1},x_i]$ se toma $x_i^*$ y se aproxima la
porción de la región por un rectángulo de altura $f(x_i^*)-g(x_i^*)$ y
base $\Delta x$. La suma de Riemann asociada es
\[
  S_n = \sum_{i=1}^n\bigl[f(x_i^*)-g(x_i^*)\bigr]\Delta x.
\]
Tomando el límite cuando $n\to\infty$ y usando la continuidad de $f-g$:
\[
  A = \lim_{n\to\infty} S_n
    = \int_a^b\bigl[f(x)-g(x)\bigr]\,dx.\qedhere
\]
\end{proof}

\begin{example}
\label{ex:parabola-recta}
Calcule el área de la región acotada por $y=x^2$ e $y=x+2$.
\begin{myproof}
\textbf{Paso 1: Análisis gráfico y región.}
Intersecciones: $x^2=x+2\Rightarrow x^2-x-2=0 \Rightarrow(x-2)(x+1)=0$,
luego $x=-1$ y $x=2$. En $[-1,2]$ se verifica $x+2\ge x^2$.
\begin{center}
% ---------------------------------------------------------------
% CORRECCIÓN: se nombran las dos trayectorias de la región
%   [a,b] = [-1, 2] y se rellena entre ellas. Luego se redibujan
%   las curvas completas (fuera de [a,b]) encima del relleno.
% ---------------------------------------------------------------
\begin{tikzpicture}[scale=0.9]
\begin{axis}[
  axis lines = middle,
  xmin = -1.8, xmax = 2.8,
  ymin = -0.6, ymax = 5.4,
  xlabel = {$x$}, ylabel = {$y$},
  grid  = major,
  grid style = {line width=0.3pt, draw=gray!25},
  tick label style = {font=\small}
]
%--- trayectorias de relleno (sólo en [-1, 2]) ---
\addplot[name path=sup, domain=-1:2, samples= 40, draw=none] {x+2};
\addplot[name path=inf, domain=-1:2, samples=100, draw=none] {x^2};
%--- región entre las dos curvas ---
\addplot[fill=violet!22, opacity=0.80] fill between[of=sup and inf];
%--- curvas completas encima del relleno ---
\addplot[domain=-1.8:2.8, samples=200, blue!80!black, thick] {x^2}
        node[pos=0.92, above left, font=\small] {$y=x^2$};
\addplot[domain=-1.8:2.8, samples=  2, red!80!black,  thick] {x+2}
        node[pos=0.85, below right, font=\small] {$y=x+2$};
%--- puntos de intersección ---
\addplot[only marks, mark=*, mark options={fill=black, scale=1.1}]
        coordinates {(-1,1)(2,4)};
\node[font=\small, above left]  at (axis cs:-1,1) {$(-1,\,1)$};
\node[font=\small, above left]  at (axis cs: 2,4) {$(2,\,4)$};
\end{axis}
\end{tikzpicture}
\end{center}
\textbf{Paso 2: Elemento diferencial.}
$dA = [\text{curva superior} - \text{curva inferior}]\,dx
     = \bigl[(x+2)-x^2\bigr]\,dx$.

\textbf{Paso 3: Límites.} $a=-1$, $b=2$.

\textbf{Paso 4: Evaluación.}
\begin{align*}
  A &= \int_{-1}^{2}\left[(x+2)-x^2\right]\,dx
     = \int_{-1}^{2}(-x^2+x+2)\,dx \\
    &= \left[-\frac{x^3}{3}+\frac{x^2}{2}+2x\right]_{-1}^{2} \\
    &= \left(-\frac{8}{3}+2+4\right) - \left(\frac{1}{3}+\frac{1}{2}-2\right) \\
    &= \frac{10}{3}+\frac{7}{6} = \frac{27}{6} = \boxed{\dfrac{9}{2}}.
\end{align*}
\end{myproof}
\end{example}

\begin{rem}
Cuando las curvas se intersectan en más de dos puntos, el área total se
obtiene sumando las integrales en cada subintervalo donde una función
domina a la otra. Es crucial graficar primero para identificar correctamente
los límites y el orden de resta en cada tramo.
\end{rem}

\begin{example}
Halle el área de la región acotada por $f(x)=\sin x$ y $g(x)=\cos x$ en
$[0,2\pi]$.

\begin{myproof}
\textbf{Paso 1: Análisis gráfico y región.}
En $[0,2\pi]$ las curvas se cruzan en $x=\pi/4$ y $x=5\pi/4$.
Analizando signos:
\begin{itemize}
  \item En $\bigl[0,\tfrac{\pi}{4}\bigr]$: $\cos x\geq\sin x$.
  \item En $\bigl[\tfrac{\pi}{4},\tfrac{5\pi}{4}\bigr]$: $\sin x\geq\cos x$.
  \item En $\bigl[\tfrac{5\pi}{4},2\pi\bigr]$: $\cos x\geq\sin x$.
\end{itemize}
\begin{center}
% ---------------------------------------------------------------
% CORRECCIÓN: tres pares de trayectorias (una por subintervalo)
%   con fillbetween. El mismo color en los tres tramos deja claro
%   que todos contribuyen positivamente al área total.
%   Se eliminan los \closedcycle que producían solapamientos.
% ---------------------------------------------------------------
\begin{tikzpicture}[scale=0.9]
\begin{axis}[
  axis lines = middle,
  xmin = -0.3, xmax = 6.6,
  ymin = -1.35, ymax = 1.35,
  xlabel = {$x$}, ylabel = {$y$},
  xtick      = {0, 0.7854, 3.1416, 3.9270, 6.2832},
  xticklabels= {$0$, $\tfrac{\pi}{4}$, $\pi$,
                $\tfrac{5\pi}{4}$, $2\pi$},
  ytick = {-1, -0.5, 0, 0.5, 1},
  grid  = major,
  grid style = {line width=0.3pt, draw=gray!25},
  tick label style = {font=\small}
]
%--- Tramo 1: [0, pi/4]  –  cos >= sin  (relleno naranja) ---
\addplot[name path=c1, domain=0:0.7854,      samples=30, draw=none]
        {cos(deg(x))};
\addplot[name path=s1, domain=0:0.7854,      samples=30, draw=none]
        {sin(deg(x))};
\addplot[fill=orange!30, opacity=0.80] fill between[of=c1 and s1];
%--- Tramo 2: [pi/4, 5pi/4]  –  sin >= cos  (relleno naranja) ---
\addplot[name path=s2, domain=0.7854:3.9270, samples=80, draw=none]
        {sin(deg(x))};
\addplot[name path=c2, domain=0.7854:3.9270, samples=80, draw=none]
        {cos(deg(x))};
\addplot[fill=orange!30, opacity=0.80] fill between[of=s2 and c2];
%--- Tramo 3: [5pi/4, 2pi]  –  cos >= sin  (relleno naranja) ---
\addplot[name path=c3, domain=3.9270:6.2832, samples=50, draw=none]
        {cos(deg(x))};
\addplot[name path=s3, domain=3.9270:6.2832, samples=50, draw=none]
        {sin(deg(x))};
\addplot[fill=orange!30, opacity=0.80] fill between[of=c3 and s3];
%--- curvas encima del relleno ---
\addplot[domain=0:6.2832, samples=300, blue!80!black, thick]
        {sin(deg(x))}
        node[pos=0.28, above, font=\small] {$\sin x$};
\addplot[domain=0:6.2832, samples=300, red!80!black,  thick]
        {cos(deg(x))}
        node[pos=0.10, above right, font=\small] {$\cos x$};
%--- puntos de cruce ---
\addplot[only marks, mark=*, mark options={fill=black, scale=0.9}]
        coordinates {(0.7854,0.7071)(3.9270,-0.7071)};
\end{axis}
\end{tikzpicture}
\end{center}
\textbf{Paso 2: Elemento diferencial.}
$dA = |\sin x - \cos x|\,dx$, dividido en tres tramos.

\textbf{Paso 3: Límites.}
$\bigl[0,\tfrac\pi4\bigr]$, $\bigl[\tfrac\pi4,\tfrac{5\pi}4\bigr]$,
$\bigl[\tfrac{5\pi}4,2\pi\bigr]$.

\textbf{Paso 4: Evaluación.}
\begin{align*}
  A &= \int_0^{\pi/4}(\cos x-\sin x)\,dx
     + \int_{\pi/4}^{5\pi/4}(\sin x-\cos x)\,dx
     + \int_{5\pi/4}^{2\pi}(\cos x-\sin x)\,dx \\
    &= \bigl[\sin x+\cos x\bigr]_0^{\pi/4}
     + \bigl[-\cos x-\sin x\bigr]_{\pi/4}^{5\pi/4}
     + \bigl[\sin x+\cos x\bigr]_{5\pi/4}^{2\pi} \\
    &= \!\left(\tfrac{\sqrt2}{2}+\tfrac{\sqrt2}{2}-(0+1)\right)
      +\!\left(\tfrac{\sqrt2}{2}+\tfrac{\sqrt2}{2}
          -\!\left(-\tfrac{\sqrt2}{2}-\tfrac{\sqrt2}{2}\right)\right)
      +\!\left((0+1)-\!\left(-\tfrac{\sqrt2}{2}-\tfrac{\sqrt2}{2}\right)\right)\\
    &= (\sqrt{2}-1)+2\sqrt{2}+(\sqrt{2}+1) = \boxed{4\sqrt{2}}.
\end{align*}
\end{myproof}
\end{example}

\begin{definition}
Si la región está acotada por curvas de la forma $x=p(y)$ y $x=q(y)$,
con $p(y)\geq q(y)$ para $y\in[c,d]$, entonces el área es
\[
  A = \int_c^d\bigl[p(y)-q(y)\bigr]\,dy.
\]
\end{definition}

\begin{rem}
Se prefiere integrar respecto a $y$ cuando las curvas se expresan
naturalmente como $x=h(y)$ (por ejemplo, parábolas horizontales) o cuando
la integración respecto a $x$ obligaría a dividir el intervalo en varios
tramos con distintas funciones superiores e inferiores.
\end{rem}

\begin{example}[Área con integración respecto a $y$]
Halle el área acotada por $x=y^2$ y $x=y+2$.

\begin{myproof}
\textbf{Paso 1: Análisis gráfico y región.}
Intersecciones: $y^2=y+2\Rightarrow y=-1,\,2$.
En $[-1,2]$ se tiene $y+2\geq y^2$ (la recta queda a la derecha).
\begin{center}
% ---------------------------------------------------------------
% CORRECCIÓN: curvas paramétricas x=h(y) trazadas con variable=\y.
%   fillbetween rellena entre la trayectoria derecha (x=y+2) y la
%   izquierda (x=y²) sin solapamientos.
% ---------------------------------------------------------------
\begin{tikzpicture}[scale=0.9]
\begin{axis}[
  axis lines = middle,
  xmin = -0.4, xmax = 4.6,
  ymin = -1.6, ymax = 2.6,
  xlabel = {$x$}, ylabel = {$y$},
  grid  = major,
  grid style = {line width=0.3pt, draw=gray!25},
  tick label style = {font=\small}
]
%--- trayectorias paramétricas invisibles para el relleno ---
\addplot[name path=der,
         variable=\y, domain=-1:2, samples=60, draw=none]
        ({\y+2}, {\y});                   % x = y+2  (derecha)
\addplot[name path=izq,
         variable=\y, domain=-1:2, samples=120, draw=none]
        ({\y^2}, {\y});                   % x = y²   (izquierda)
%--- región entre las dos curvas ---
\addplot[fill=violet!22, opacity=0.80] fill between[of=der and izq];
%--- curvas completas encima del relleno ---
\addplot[variable=\y, domain=-1.6:2.6, samples=200,
         blue!80!black, thick]
        ({\y^2}, {\y})
        node[pos=0.82, right, font=\small] {$x=y^2$};
\addplot[variable=\y, domain=-1.6:2.6, samples=2,
         red!80!black, thick]
        ({\y+2}, {\y})
        node[pos=0.85, below right, font=\small] {$x=y+2$};
%--- puntos de intersección ---
\addplot[only marks, mark=*, mark options={fill=black, scale=1.1}]
        coordinates {(1,-1)(4,2)};
\node[font=\small, left]  at (axis cs:1,-1) {$(-1)$};
\node[font=\small, above left] at (axis cs:4, 2) {$(2)$};
\end{axis}
\end{tikzpicture}
\end{center}
\textbf{Paso 2: Elemento diferencial.}
$dA = [\text{curva derecha} - \text{curva izquierda}]\,dy
     = \bigl[(y+2)-y^2\bigr]\,dy$.

\textbf{Paso 3: Límites.} $y\in[-1,2]$.

\textbf{Paso 4: Evaluación.}
\[
  A = \int_{-1}^{2}\bigl[(y+2)-y^2\bigr]\,dy
    = \left[\frac{y^2}{2}+2y-\frac{y^3}{3}\right]_{-1}^{2}
    = \left(2+4-\frac{8}{3}\right) - \left(\frac{1}{2}-2+\frac{1}{3}\right)
    = \boxed{\dfrac{9}{2}},
\]
coincidiendo con el Ejemplo~\ref{ex:parabola-recta}, como debe ser.
\end{myproof}
\end{example}

\begin{example}
Halle el área entre $y=\sqrt{x}$, $y=x-2$ y el eje $y$ en el primer
cuadrante.

\begin{myproof}
\textbf{Paso 1: Análisis gráfico y región.}
Las curvas se cortan cuando $\sqrt{x}=x-2 \Rightarrow x=4 \Rightarrow y=2$.
Despejando: $x=y^2$ (parábola, izquierda) y $x=y+2$ (recta, derecha),
con $y\in[0,2]$.
\begin{center}
% ---------------------------------------------------------------
% CORRECCIÓN: misma técnica paramétrica variable=\y.
%   El eje y (x=0) actúa de límite izquierdo implícito: como
%   y+2 > y² para y∈[0,2], basta con la región entre las dos
%   curvas sobre ese intervalo.
% ---------------------------------------------------------------
\begin{tikzpicture}[scale=0.9]
\begin{axis}[
  axis lines = middle,
  xmin = -0.3, xmax = 4.8,
  ymin = -0.3, ymax = 2.6,
  xlabel = {$x$}, ylabel = {$y$},
  grid  = major,
  grid style = {line width=0.3pt, draw=gray!25},
  tick label style = {font=\small}
]
%--- trayectorias invisibles ---
\addplot[name path=der,
         variable=\y, domain=0:2, samples=50, draw=none]
        ({\y+2}, {\y});
\addplot[name path=izq,
         variable=\y, domain=0:2, samples=100, draw=none]
        ({\y^2}, {\y});
%--- relleno ---
\addplot[fill=violet!22, opacity=0.80] fill between[of=der and izq];
%--- curvas completas ---
\addplot[variable=\y, domain=0:2.5, samples=200, blue!80!black, thick]
        ({\y^2}, {\y})
        node[pos=0.75, right, font=\small] {$x=y^2$};
\addplot[variable=\y, domain=0:2.5, samples=2,   red!80!black,  thick]
        ({\y+2}, {\y})
        node[pos=0.80, below right, font=\small] {$x=y+2$};
%--- punto de intersección ---
\addplot[only marks, mark=*, mark options={fill=black, scale=1.1}]
        coordinates {(4,2)};
\node[font=\small, above left] at (axis cs:4,2) {$(4,\,2)$};
\end{axis}
\end{tikzpicture}
\end{center}
\textbf{Paso 2: Elemento diferencial.}
$dA = \bigl[(y+2)-y^2\bigr]\,dy$.

\textbf{Paso 3: Límites.} $y\in[0,2]$.

\textbf{Paso 4: Evaluación.}
\[
  A = \int_0^2\bigl[(y+2)-y^2\bigr]\,dy
    = \left[\frac{y^2}{2}+2y-\frac{y^3}{3}\right]_0^2
    = 2+4-\frac{8}{3} = \boxed{\dfrac{10}{3}}.
\]
\end{myproof}
\end{example}

% ============================================================
\section{Volúmenes de sólidos de revolución}
% ============================================================

Un \textbf{sólido de revolución} se genera al rotar una región plana
alrededor de una recta llamada \emph{eje de revolución}. Presentamos tres
métodos fundamentales, cada uno adecuado para distintas configuraciones.

% -----------------------------------------------------------
\subsection{Método de los discos}
% -----------------------------------------------------------

\begin{theorem}[Volumen por discos respecto al eje $x$]
Sea $R$ la región acotada por $y=f(x)\geq 0$, $x=a$, $x=b$ y el eje $x$.
Al rotar $R$ alrededor del eje $x$, el volumen del sólido generado es
\[
  V = \pi\int_a^b\bigl[f(x)\bigr]^2\,dx.
\]
\end{theorem}

\begin{proof}
Dividimos $[a,b]$ en $n$ subintervalos de ancho $\Delta x$. La porción del
sólido sobre $[x_{i-1},x_i]$ se aproxima por un disco circular de radio
$f(x_i^*)$ y grosor $\Delta x$, cuyo volumen es $\pi[f(x_i^*)]^2\Delta x$.
Sumando y tomando límite:
\[
  V = \lim_{n\to\infty}\sum_{i=1}^n\pi[f(x_i^*)]^2\Delta x
    = \pi\int_a^b[f(x)]^2\,dx.\qedhere
\]
\end{proof}

\begin{example}
Demuestre que el volumen de un cono de radio $R$ y altura $h$ es
$V=\frac{1}{3}\pi R^2 h$.

\begin{myproof}
\textbf{Paso 1: Generatriz y región.} La recta que genera el cono es
$y=\frac{R}{h}x$, $x\in[0,h]$.

\textbf{Paso 2: Elemento diferencial.} Disco de radio $f(x)$ y grosor $dx$:
$dV = \pi [f(x)]^2 dx$.

\textbf{Paso 3: Límites.} $x$ va de $0$ a $h$.

\textbf{Paso 4: Evaluación.}
\[
  V = \pi\int_0^h\!\left(\frac{R}{h}x\right)^{\!2}dx
    = \frac{\pi R^2}{h^2}\int_0^h x^2\,dx
    = \frac{\pi R^2}{h^2}\cdot\frac{h^3}{3}
    = \boxed{\frac{1}{3}\pi R^2 h}.
\]
\end{myproof}
\end{example}

\begin{example}
Demuestre que el volumen de una esfera de radio $R$ es
$V = \frac{4}{3}\pi R^3$.

\begin{myproof}
\textbf{Paso 1: Generatriz y región.} Rotamos la semicircunferencia
$y=\sqrt{R^2-x^2}$, $x\in[-R,R]$, alrededor del eje $x$.

\textbf{Paso 2: Elemento diferencial.}
$dV = \pi y^2 dx = \pi(R^2-x^2)dx$.

\textbf{Paso 3: Límites.} $x\in[-R,R]$. Por simetría par,
$V = 2\pi\int_0^R (R^2-x^2)dx$.

\textbf{Paso 4: Evaluación.}
\[
  V = 2\pi\left[R^2 x-\frac{x^3}{3}\right]_0^R
    = 2\pi\left(R^3-\frac{R^3}{3}\right)
    = 2\pi\cdot\frac{2R^3}{3}
    = \boxed{\dfrac{4}{3}\pi R^3}.\qedhere
\]
\end{myproof}
\end{example}

% -----------------------------------------------------------
\subsection{Método de las arandelas (\emph{washers})}
% -----------------------------------------------------------

\begin{theorem}[Volumen por arandelas]
Sea $R$ la región acotada por $y=f(x)$ (superior) e $y=g(x)$ (inferior),
con $f(x)\geq g(x)\geq 0$, $x\in[a,b]$. Al rotar $R$ alrededor del eje $x$,
el volumen es
\[
  V = \pi\int_a^b\Bigl(\bigl[f(x)\bigr]^2-\bigl[g(x)\bigr]^2\Bigr)\,dx.
\]
\end{theorem}

\begin{rem}
La arandela en la posición $x$ tiene radio externo $R(x)=f(x)$ y radio
interno $r(x)=g(x)$. Su área transversal es
$\pi R(x)^2-\pi r(x)^2$, de donde se obtiene la fórmula integrando sobre
$[a,b]$.
\end{rem}

\begin{example}
La región entre $y=\sqrt{x}$ e $y=x^2$ en $[0,1]$ gira alrededor del eje
$x$. Calcule el volumen.

\begin{myproof}
\textbf{Paso 1: Región y eje.}
En $[0,1]$: $\sqrt{x}\geq x^2$.
Radio externo $R(x)=\sqrt{x}$, radio interno $r(x)=x^2$.
\begin{center}
% ---------------------------------------------------------------
% CORRECCIÓN: fillbetween entre la curva superior y la inferior.
%   Se anotan los radios para reforzar el concepto de arandela.
% ---------------------------------------------------------------
\begin{tikzpicture}[scale=0.9]
\begin{axis}[
  axis lines = middle,
  xmin = -0.1, xmax = 1.15,
  ymin = -0.1, ymax = 1.15,
  xlabel = {$x$}, ylabel = {$y$},
  xtick = {0, 0.25, 0.5, 0.75, 1},
  ytick = {0, 0.25, 0.5, 0.75, 1},
  grid  = major,
  grid style = {line width=0.3pt, draw=gray!25},
  tick label style = {font=\small}
]
%--- trayectorias invisibles ---
\addplot[name path=sup, domain=0:1, samples=120, draw=none] {sqrt(x)};
\addplot[name path=inf, domain=0:1, samples=120, draw=none] {x^2};
%--- región entre curvas ---
\addplot[fill=violet!22, opacity=0.80] fill between[of=sup and inf];
%--- curvas encima del relleno ---
\addplot[domain=0:1.1, samples=200, blue!80!black, thick] {sqrt(x)}
        node[pos=0.80, above left, font=\small] {$y=\sqrt{x}$};
\addplot[domain=0:1.1, samples=200, red!80!black,  thick] {x^2}
        node[pos=0.85, below right, font=\small] {$y=x^2$};
%--- radio externo e interno en x=0.5 ---
\draw[<->, thick, blue!60!black]
     (axis cs:0.5,0) -- (axis cs:0.5,{sqrt(0.5)})
     node[midway, right, font=\scriptsize] {$R(x)$};
\draw[<->, thick, red!60!black]
     (axis cs:0.5,0) -- (axis cs:0.5,{0.5^2})
     node[midway, left, font=\scriptsize] {$r(x)$};
\end{axis}
\end{tikzpicture}
\end{center}
\textbf{Paso 2: Elemento diferencial.}
Arandela: $dV = \pi\bigl[R(x)^2 - r(x)^2\bigr]dx = \pi\bigl[x - x^4\bigr]dx$.

\textbf{Paso 3: Límites.} $x\in[0,1]$.

\textbf{Paso 4: Evaluación.}
\begin{align*}
  V &= \pi\int_0^1(x-x^4)\,dx
     = \pi\left[\frac{x^2}{2}-\frac{x^5}{5}\right]_0^1
     = \pi\left(\frac{1}{2}-\frac{1}{5}\right)
     = \boxed{\dfrac{3\pi}{10}}.
\end{align*}
\end{myproof}
\end{example}

\begin{example}
La región acotada por $y=x^2$ e $y=1$ gira alrededor de la recta $y=2$.
Calcule el volumen.

\begin{myproof}
\textbf{Paso 1: Región y eje.}
Intersecciones: $x^2=1\Rightarrow x=\pm 1$.
El eje de rotación $y=2$ queda por encima de toda la región.
\begin{center}
% ---------------------------------------------------------------
% CORRECCIÓN: la región de integración (entre y=x² e y=1) se
%   rellena limpiamente con fillbetween. El eje de rotación y=2
%   se traza como línea punteada con etiqueta. Los radios se
%   anotan en x=0.6 para reforzar el concepto de arandela.
% ---------------------------------------------------------------
\begin{tikzpicture}[scale=0.9]
\begin{axis}[
  axis lines = middle,
  xmin = -1.6, xmax = 1.6,
  ymin = -0.2, ymax = 2.4,
  xlabel = {$x$}, ylabel = {$y$},
  xtick = {-1, -0.5, 0, 0.5, 1},
  ytick = {0, 0.5, 1, 1.5, 2},
  grid  = major,
  grid style = {line width=0.3pt, draw=gray!25},
  tick label style = {font=\small}
]
%--- eje de rotación y=2 ---
\draw[dashed, gray!60, thick]
     (axis cs:-1.6,2) -- (axis cs:1.6,2)
     node[right, font=\small, gray!70!black] {eje $y=2$};
%--- trayectorias invisibles en [-1,1] ---
\addplot[name path=Rsup, domain=-1:1, samples=2,   draw=none] {1};
\addplot[name path=Rinf, domain=-1:1, samples=100, draw=none] {x^2};
%--- región entre y=x² e y=1 ---
\addplot[fill=violet!22, opacity=0.80] fill between[of=Rsup and Rinf];
%--- curvas encima del relleno ---
\addplot[domain=-1.5:1.5, samples=200, blue!80!black, thick] {x^2}
        node[pos=0.93, above left, font=\small] {$y=x^2$};
\draw[red!80!black, thick]
     (axis cs:-1.5,1) -- (axis cs:1.5,1)
     node[right, font=\small] {$y=1$};
%--- radios en x=0.6 ---
\draw[<->, thick, blue!60!black]
     (axis cs:0.6,{0.6^2}) -- (axis cs:0.6,2)
     node[midway, right, font=\scriptsize] {$R(x)$};
\draw[<->, thick, red!60!black]
     (axis cs:-0.6,1) -- (axis cs:-0.6,2)
     node[midway, left, font=\scriptsize] {$r(x)$};
%--- puntos de intersección ---
\addplot[only marks, mark=*, mark options={fill=black, scale=0.9}]
        coordinates {(-1,1)(1,1)};
\end{axis}
\end{tikzpicture}
\end{center}
\textbf{Paso 2: Radios y elemento diferencial.}
Radio externo (de $y=2$ a $y=x^2$): $R(x)=2-x^2$.
Radio interno (de $y=2$ a $y=1$): $r(x)=2-1=1$.
\[
  dV = \pi\bigl[R(x)^2 - r(x)^2\bigr]dx
     = \pi\bigl[(2-x^2)^2 - 1\bigr]dx.
\]
\textbf{Paso 3: Límites.} $x\in[-1,1]$. Por simetría,
$V = 2\pi\int_0^1 \bigl[(2-x^2)^2-1\bigr]dx$.

\textbf{Paso 4: Evaluación.}
\begin{align*}
  V &= \pi\int_{-1}^{1}\left[(2-x^2)^2-1\right]\,dx
     = 2\pi\int_0^1(3-4x^2+x^4)\,dx \\
    &= 2\pi\left[3x-\frac{4x^3}{3}+\frac{x^5}{5}\right]_0^1
     = 2\pi\left(3-\frac{4}{3}+\frac{1}{5}\right)
     = \boxed{\dfrac{56\pi}{15}}.
\end{align*}
\end{myproof}
\end{example}

% -----------------------------------------------------------
\subsection{Método de los cascarones cilíndricos (\emph{shells})}
% -----------------------------------------------------------

\begin{theorem}[Volumen por cascarones girando respecto al eje $y$]
Sea $R$ la región acotada por $y=f(x)\geq 0$, $x=a$, $x=b$
($0\leq a<b$) y el eje $x$. Al rotar $R$ alrededor del eje $y$, el volumen
es
\[
  V = 2\pi\int_a^b x\,f(x)\,dx,
\]
donde $x$ es el radio promedio del cascarón y $f(x)$ su altura.
\end{theorem}

\begin{proof}
Un rectángulo vertical de ancho $\Delta x$, altura $f(x_i^*)$ y distancia
$x_i^*$ al eje $y$ genera, al rotar, un cascarón cilíndrico cuyo volumen
es aproximadamente
\[
  2\pi x_i^*\cdot f(x_i^*)\cdot\Delta x
  \quad\bigl(\text{circunferencia}\times\text{altura}\times\text{grosor}\bigr).
\]
Sumando y tomando límite se obtiene la fórmula.
\end{proof}

\begin{rem}
El método de cascarones es especialmente útil cuando:
\begin{itemize}
  \item El eje de revolución es paralelo al rectángulo representativo.
  \item Despejar $x$ en términos de $y$ es algebraicamente complejo.
  \item El sólido tiene huecos que complican el método de arandelas.
\end{itemize}
Como regla práctica: si el eje de rotación es el eje $y$ (o una recta
vertical), pruebe primero con cascarones integrando respecto a $x$.
Si el eje de rotación es el eje $x$ (o una recta horizontal), pruebe
primero con discos/arandelas.
\end{rem}

\begin{example}
Calcule el volumen generado al rotar la región bajo $y=\sqrt{x}$ en $[0,4]$
alrededor del eje $y$.

\begin{myproof}
\begin{center}
% ---------------------------------------------------------------
% CORRECCIÓN: se dibuja la región bajo y=√x con fillbetween y se
%   muestra un cascarón representativo (rectángulo vertical) para
%   ilustrar el método, en lugar de una flecha circular ambigua.
% ---------------------------------------------------------------
\begin{tikzpicture}[scale=0.9]
\begin{axis}[
  axis lines = middle,
  xmin = -0.3, xmax = 4.6,
  ymin = -0.3, ymax = 2.4,
  xlabel = {$x$}, ylabel = {$y$},
  xtick = {0, 1, 2, 3, 4},
  ytick = {0, 0.5, 1, 1.5, 2},
  grid  = major,
  grid style = {line width=0.3pt, draw=gray!25},
  tick label style = {font=\small}
]
%--- trayectorias invisibles ---
\addplot[name path=cur,  domain=0:4, samples=200, draw=none] {sqrt(x)};
\addplot[name path=base, domain=0:4, samples=  2, draw=none] {0};
%--- región bajo la curva ---
\addplot[fill=violet!22, opacity=0.80] fill between[of=cur and base];
%--- cascarón representativo en x=2.5, ancho dx=0.3 ---
\addplot[fill=orange!40, opacity=0.70, draw=none]
        coordinates {(2.5,0)(2.8,0)(2.8,{sqrt(2.8)})(2.5,{sqrt(2.5)})(2.5,0)};
\draw[orange!70!black, thick]
     (axis cs:2.5,0) rectangle (axis cs:2.8,{sqrt(2.65)});
%--- anotación de radio y altura ---
\draw[<->, thick, gray!70!black]
     (axis cs:0,-0.15) -- (axis cs:2.65,-0.15)
     node[midway, below, font=\scriptsize] {$x$};
\draw[<->, thick, blue!60!black]
     (axis cs:4.35,0) -- (axis cs:4.35,{sqrt(4)})
     node[midway, right, font=\scriptsize] {$\sqrt{x}$};
%--- curva encima ---
\addplot[domain=0:4.5, samples=200, blue!80!black, thick] {sqrt(x)}
        node[pos=0.95, above left, font=\small] {$y=\sqrt{x}$};
\end{axis}
\end{tikzpicture}
\end{center}
\textbf{Solución por cascarones:}

\textbf{Paso 1: Elemento diferencial.}
Cascarón vertical: $dV = 2\pi\,(\text{radio})\,(\text{altura})\,dx
= 2\pi x \sqrt{x}\,dx = 2\pi x^{3/2}dx$.

\textbf{Paso 2: Límites.} $x\in[0,4]$.

\textbf{Paso 3: Evaluación.}
\[
  V = 2\pi\int_0^4 x^{3/2}\,dx
    = 2\pi\left[\frac{2}{5}x^{5/2}\right]_0^4
    = \frac{4\pi}{5}\cdot 32
    = \boxed{\dfrac{128\pi}{5}}.
\]

\textbf{Verificación por arandelas respecto a $y$:}
Despejamos $x=y^2$, con $y\in[0,2]$. El sólido tiene radio externo $4$ y
radio interno $y^2$:
\[
  V = \pi\int_0^2\bigl(4^2-(y^2)^2\bigr)\,dy
    = \pi\int_0^2(16-y^4)\,dy
    = \pi\left[16y-\frac{y^5}{5}\right]_0^2
    = \pi\!\left(32-\frac{32}{5}\right)
    = \frac{128\pi}{5}.\qedhere
\]
\end{myproof}
\end{example}

\begin{example}
La región acotada por $y=4x-x^2$ e $y=0$ gira alrededor de la recta
$x=5$. Calcule el volumen.
\begin{myproof}
\textbf{Paso 1: Región y eje.}
La parábola toca el eje $x$ en $x=0$ y $x=4$. El eje de rotación es $x=5$.

\textbf{Paso 2: Elemento diferencial.}
Cascarón cilíndrico. Radio $= 5-x$. Altura $= 4x-x^2$.
\[
  dV = 2\pi (5-x)(4x-x^2)\,dx.
\]
\textbf{Paso 3: Límites.} $x\in[0,4]$.

\textbf{Paso 4: Evaluación.}
\begin{align*}
  V &= 2\pi\int_0^4(5-x)(4x-x^2)\,dx
     = 2\pi\int_0^4(20x-9x^2+x^3)\,dx \\
    &= 2\pi\left[10x^2-3x^3+\frac{x^4}{4}\right]_0^4
     = 2\pi(160-192+64)
     = \boxed{64\pi}.
\end{align*}
\end{myproof}
\end{example}

% ============================================================
\section{Longitud de arco}
% ============================================================

\begin{definition}[Longitud de una curva suave]
Sea $f$ una función con derivada continua en $[a,b]$. La
\textbf{longitud de arco} de la gráfica de $f$ desde $x=a$ hasta $x=b$ es
\[
  s = \int_a^b\sqrt{1+\bigl[f'(x)\bigr]^2}\,dx.
\]
\end{definition}

\begin{proof}[Esquema de demostración]
Dividimos $[a,b]$ en $n$ subintervalos. La longitud del segmento que une
$(x_{i-1},f(x_{i-1}))$ con $(x_i,f(x_i))$ es, por el teorema de Pitágoras,
\[
  \Delta s_i = \sqrt{(\Delta x_i)^2+(\Delta y_i)^2}
             = \sqrt{1+\!\left(\frac{\Delta y_i}{\Delta x_i}\right)^{\!2}}
               \,\Delta x_i.
\]
Por el Teorema del Valor Medio, $\dfrac{\Delta y_i}{\Delta x_i}=f'(c_i)$
para algún $c_i\in(x_{i-1},x_i)$. Sumando y tomando límite:
\[
  s = \lim_{n\to\infty}\sum_{i=1}^n\sqrt{1+[f'(c_i)]^2}\,\Delta x_i
    = \int_a^b\sqrt{1+[f'(x)]^2}\,dx.\qedhere
\]
\end{proof}

\begin{rem}
El elemento diferencial de longitud de arco es
$ds=\sqrt{1+[f'(x)]^2}\,dx = \sqrt{dx^2+dy^2}$, que corresponde
geométricamente a la hipotenusa del triángulo diferencial de catetos $dx$
y $dy=f'(x)\,dx$. Esta interpretación facilita recordar y derivar la
fórmula en distintos contextos (paramétrico, polar, etc.).
\end{rem}

\begin{example}[Longitud de una parábola]
Calcule la longitud de arco de $y=x^2$ en $[0,1]$.

\begin{myproof}
\textbf{Paso 1: Derivada e integrando.}
$f'(x)=2x \Rightarrow ds = \sqrt{1+4x^2}\,dx$.

\textbf{Paso 2: Límites.} $x\in[0,1]$.

\textbf{Paso 3: Evaluación.}
Con la sustitución trigonométrica $2x=\tan\theta$,
$dx=\frac{1}{2}\sec^2\theta\,d\theta$:
\begin{align*}
  s &= \frac{1}{2}\int_0^{\arctan 2}\sec^3\theta\,d\theta
     = \frac{1}{4}\Bigl[\sec\theta\tan\theta
       +\ln|\sec\theta+\tan\theta|\Bigr]_0^{\arctan 2} \\
    &= \frac{1}{4}\bigl(\sqrt{5}\cdot 2+\ln(2+\sqrt{5})\bigr)
     = \boxed{\dfrac{\sqrt{5}}{2}+\dfrac{1}{4}\ln(2+\sqrt{5})\approx 1.479}.
\end{align*}
\end{myproof}
\end{example}

\begin{example}[Longitud del arco de la catenaria]
Calcule la longitud de arco de $y=\cosh x = \dfrac{e^x+e^{-x}}{2}$ en
$[0,\ln 2]$.

\begin{myproof}
\textbf{Paso 1: Derivada e integrando.}
$f'(x)=\sinh x$. Usando $\cosh^2 x-\sinh^2 x=1$:
$ds = \sqrt{1+\sinh^2 x}\,dx = \cosh x\,dx$.

\textbf{Paso 2: Límites.} $x\in[0,\ln 2]$.

\textbf{Paso 3: Evaluación.}
\[
  s = \int_0^{\ln 2}\cosh x\,dx
    = \bigl[\sinh x\bigr]_0^{\ln 2}
    = \sinh(\ln 2)
    = \frac{e^{\ln 2} - e^{-\ln 2}}{2}
    = \frac{2 - \frac{1}{2}}{2}
    = \boxed{\dfrac{3}{4}}.
\]
\end{myproof}
\end{example}

% ============================================================
\section{Áreas de superficies de revolución}
% ============================================================

\begin{definition}[Área superficial de revolución --- eje $x$]
Sea $f$ una función no negativa con derivada continua en $[a,b]$. El área
de la superficie generada al rotar la gráfica de $f$ alrededor del eje $x$
es
\[
  S = 2\pi\int_a^b f(x)\sqrt{1+[f'(x)]^2}\,dx.
\]
\end{definition}

\begin{proof}[Idea de demostración]
Un pequeño arco de longitud $ds=\sqrt{1+[f'(x)]^2}\,dx$, al rotar alrededor
del eje $x$, genera una banda (frustum de cono) cuya área aproximada es
\[
  dS = 2\pi f(x)\,ds = 2\pi f(x)\sqrt{1+[f'(x)]^2}\,dx.
\]
Integrando: $S = 2\pi\int_a^b f(x)\sqrt{1+[f'(x)]^2}\,dx$. \qedhere
\end{proof}

\begin{rem}
El factor $2\pi f(x)$ es la circunferencia del círculo descrito por el
punto $(x,f(x))$ al rotar. El factor $ds$ es el elemento de longitud de
arco, que ``estira'' la banda según la inclinación de la curva. Si se
tomara $ds=dx$ en lugar de $ds=\sqrt{1+[f'(x)]^2}\,dx$ se obtendría la
fórmula incorrecta para un cilindro, no para la superficie real.
\end{rem}

\begin{example}[Superficie de una esfera]
Demuestre que el área de una esfera de radio $R$ es $4\pi R^2$.

\begin{myproof}
\textbf{Paso 1: Generatriz y derivada.}
Rotamos $y=\sqrt{R^2-x^2}$, $x\in[-R,R]$, alrededor del eje $x$:
\[
  y' = \frac{-x}{\sqrt{R^2-x^2}} \implies
  1+(y')^2 = 1 + \frac{x^2}{R^2-x^2} = \frac{R^2}{R^2-x^2}.
\]
\textbf{Paso 2: Elemento diferencial.}
\[
  dS = 2\pi y \sqrt{1+(y')^2}\,dx
     = 2\pi \sqrt{R^2-x^2}\cdot\frac{R}{\sqrt{R^2-x^2}}\,dx
     = 2\pi R\,dx.
\]
\textbf{Paso 3: Límites.} $x\in[-R,R]$.

\textbf{Paso 4: Evaluación.}
\[
  S = 2\pi R\int_{-R}^{R}dx
    = 2\pi R \cdot 2R
    = \boxed{4\pi R^2}.\qedhere
\]
\end{myproof}
\end{example}

\begin{example}[Superficie de un paraboloide]
Halle el área de la superficie generada al rotar $y=x^2$ en $[0,1]$
alrededor del eje $y$.

\begin{myproof}
\textbf{Paso 1: Elemento diferencial.}
Al rotar alrededor del eje $y$, el radio de cada banda es $x$.
\[
  dS = 2\pi x \sqrt{1+(2x)^2}\,dx = 2\pi x \sqrt{1+4x^2}\,dx.
\]
\textbf{Paso 2: Límites.} $x\in[0,1]$.

\textbf{Paso 3: Evaluación.}
Con $u=1+4x^2$, $du=8x\,dx$:
\[
  S = 2\pi\int_1^5\sqrt{u}\,\frac{du}{8}
    = \frac{\pi}{4}\cdot\frac{2}{3}\bigl[u^{3/2}\bigr]_1^5
    = \frac{\pi}{6}(5\sqrt{5}-1)
    = \boxed{\dfrac{\pi(5\sqrt{5}-1)}{6}}.
\]
\end{myproof}
\end{example}

\begin{theorem}[Superficie alrededor del eje $y$]
Si la gráfica de $y=f(x)$, $x\in[a,b]$, con $f\geq 0$ y $f'$ continua,
gira alrededor del eje $y$, el área superficial es
\[
  S = 2\pi\int_a^b x\sqrt{1+[f'(x)]^2}\,dx.
\]
\end{theorem}

% ============================================================
\section{Aplicaciones físicas}
% ============================================================

\subsection{Trabajo mecánico}

\begin{definition}[Trabajo de una fuerza variable]
Si una fuerza $F(x)$ actúa en la dirección del movimiento a lo largo del
eje $x$, el trabajo realizado al mover un objeto desde $x=a$ hasta $x=b$
es
\[
  W = \int_a^b F(x)\,dx.
\]
\end{definition}

\begin{rem}
Cuando la fuerza es constante, $W=F\cdot(b-a)$, que recupera la fórmula
elemental $W=Fd$. La integral generaliza este concepto a fuerzas variables.
Las unidades del trabajo son joules ($\text{N}\cdot\text{m}$) en el
Sistema Internacional.
\end{rem}

\begin{example}[Ley de Hooke para la longitud de un resorte]
Un resorte requiere $10\,\text{N}$ para estirarse $0.1\,\text{m}$ desde
su longitud natural. Calcule el trabajo para estirarlo $0.3\,\text{m}$.

\begin{myproof}
\textbf{Paso 1: Modelo de fuerza.}
Ley de Hooke: $F(x)=kx$. De $10=k(0.1)$ obtenemos $k=100\,\text{N/m}$.

\textbf{Paso 2: Elemento diferencial.} $dW = F(x)\,dx = 100x\,dx$.

\textbf{Paso 3: Límites.} Desde $x=0$ hasta $x=0.3$.

\textbf{Paso 4: Evaluación.}
\[
  W = \int_0^{0.3}100x\,dx
    = 100\left[\frac{x^2}{2}\right]_0^{0.3}
    = 50\cdot 0.09
    = \boxed{4.5\,\text{J}}.
\]
\end{myproof}
\end{example}

\begin{example}
El mismo resorte del ejemplo anterior está comprimido $0.05\,\text{m}$.
¿Cuánto trabajo adicional se requiere para comprimirlo $0.15\,\text{m}$
más?

\begin{myproof}
\textbf{Paso 1: Límites.}
El resorte ya está en $x=-0.05\,\text{m}$ y lo llevamos a
$x=-0.20\,\text{m}$. Como el trabajo depende de la magnitud del
desplazamiento y la fuerza se opone, usamos coordenadas positivas para
la compresión: de $0.05$ a $0.20$.

\textbf{Paso 2: Evaluación.}
\[
  W = \int_{0.05}^{0.20}100x\,dx
    = 50\left[x^2\right]_{0.05}^{0.20}
    = 50(0.04-0.0025)
    = \boxed{1.875\,\text{J}}.
\]
\end{myproof}
\end{example}

\begin{example}
Un tanque en forma de cono invertido tiene radio superior $3\,\text{m}$
y altura $6\,\text{m}$, y está lleno de agua
($\rho=1000\,\text{kg/m}^3$, $g=9.8\,\text{m/s}^2$). Calcule el trabajo
para bombear toda el agua hasta el borde.

\begin{myproof}
\textbf{Paso 1: Geometría y rebanada.}
Situamos el vértice en el origen. A altura $y$, por semejanza de
triángulos, el radio es $r(y)=\frac{3}{6}y=\frac{y}{2}$.

\textbf{Paso 2: Elemento diferencial.}
\begin{alignat*}{2}
  dV  &= \pi r(y)^2 dy &&= \frac{\pi y^2}{4}\,dy. \\
  dF  &= \rho g\,dV   &&= \frac{9800\pi y^2}{4}\,dy. \\
  dW  &= dF \cdot (6-y) &&= \frac{9800\pi}{4}\,y^2(6-y)\,dy.
\end{alignat*}
\textbf{Paso 3: Límites.} $y\in[0,6]$.

\textbf{Paso 4: Evaluación.}
\begin{align*}
  W &= \frac{9800\pi}{4}\int_0^6(6y^2-y^3)\,dy
     = 2450\pi\left[2y^3-\frac{y^4}{4}\right]_0^6 \\
    &= 2450\pi\left(2\cdot 216 - 324\right)
     = 2450\pi(108)
     = \boxed{264600\pi\approx 8.31\times 10^5\,\text{J}}.
\end{align*}
\end{myproof}
\end{example}

\begin{example}
Un tanque cilíndrico de radio $2\,\text{m}$ y altura $5\,\text{m}$ está
lleno de agua. Calcule el trabajo para bombear toda el agua hasta el borde.

\begin{myproof}
\textbf{Paso 1: Geometría.}
Área transversal constante $A=\pi(2)^2=4\pi$.

\textbf{Paso 2: Elemento diferencial.}
Peso de rebanada: $dF = \rho g A\,dy = 1000\cdot 9.8\cdot 4\pi\,dy
= 39200\pi\,dy$. Distancia: $5-y$.

\textbf{Paso 3: Límites.} $y\in[0,5]$.

\textbf{Paso 4: Evaluación.}
\[
  W = \int_0^5 39200\pi(5-y)\,dy
    = 39200\pi\left[5y-\frac{y^2}{2}\right]_0^5
    = 39200\pi\cdot\frac{25}{2}
    = \boxed{490000\pi\approx 1.54\times 10^6\,\text{J}}.
\]
\end{myproof}
\end{example}

% -----------------------------------------------------------
\subsection{Centro de masa y centroide}
% -----------------------------------------------------------

\begin{definition}[Momentos y centro de masa de una lámina plana]
Sea $R$ una región plana de densidad constante $\rho$, acotada por
$y=f(x)\geq 0$, $x=a$ y $x=b$. Sus \textbf{momentos} respecto a los ejes
son
\[
  M_y = \rho\int_a^b x\,f(x)\,dx, \qquad
  M_x = \frac{\rho}{2}\int_a^b [f(x)]^2\,dx.
\]
La \textbf{masa} es $m=\rho\int_a^b f(x)\,dx$ y el
\textbf{centro de masa} $(\bar x,\bar y)$ satisface
\[
  \bar{x} = \frac{M_y}{m}, \qquad \bar{y} = \frac{M_x}{m}.
\]
Cuando $\rho=1$ (o se cancela), $(\bar x,\bar y)$ se llama
\textbf{centroide} de la región.
\end{definition}

\begin{rem}
El centro de masa es el punto donde se equilibraría perfectamente la lámina
si se apoyara en él. Para regiones con densidad variable $\rho=\rho(x,y)$,
los momentos se calculan mediante integrales dobles, que se estudiarán en
Cálculo Multivariable.
\end{rem}

\begin{theorem}[Teorema de Pappus para el volumen]
Si una región plana $R$ de área $A$ se hace girar alrededor de un eje
exterior (que no la cruza), el volumen del sólido de revolución generado
es
\[
  V = 2\pi\,\bar{d}\cdot A,
\]
donde $\bar{d}$ es la distancia del centroide de $R$ al eje de rotación.
\end{theorem}

\begin{rem}
El Teorema de Pappus permite calcular volúmenes de revolución a partir del
centroide (a veces más sencillo), o bien determinar el centroide de una
región a partir de un volumen conocido.
\end{rem}

\begin{example}[Centroide de un triángulo rectángulo]
Halle el centroide de la región acotada por $y=1-x$, $x=0$ e $y=0$.

\begin{myproof}
\textbf{Paso 1: Geometría y masa.}
Tomamos $\rho=1$. Área/Masa: $m=\int_0^1(1-x)\,dx=\frac{1}{2}$.

\textbf{Paso 2: Momentos.}
\[
  M_y = \int_0^1 x(1-x)\,dx
      = \left[\frac{x^2}{2}-\frac{x^3}{3}\right]_0^1 = \frac{1}{6}, \quad
  M_x = \frac{1}{2}\int_0^1(1-x)^2\,dx = \frac{1}{2}\cdot\frac{1}{3}
      = \frac{1}{6}.
\]
\textbf{Paso 3: Coordenadas.}
$\bar{x}=\dfrac{1/6}{1/2}=\dfrac{1}{3}$,
$\bar{y}=\dfrac{1/6}{1/2}=\dfrac{1}{3}$.
\[
  \boxed{\left(\dfrac{1}{3},\dfrac{1}{3}\right)}
\]
Resultado consistente con la propiedad geométrica de que el centroide de
un triángulo se ubica en el tercio de la altura desde cada base.
\end{myproof}
\end{example}

\begin{example}[Centroide de un semicírculo]
Halle el centroide de la región semicircular $x^2+y^2\leq R^2$, $y\geq 0$.

\begin{myproof}
\textbf{Paso 1: Simetría y área.}
Por simetría respecto al eje $y$, $\bar{x}=0$.
El área es $A=\frac{\pi R^2}{2}$.

\textbf{Paso 2: Momento $M_x$.}
\[
  M_x = \frac{1}{2}\int_{-R}^{R}(R^2-x^2)\,dx
      = \frac{1}{2}\left[R^2 x-\frac{x^3}{3}\right]_{-R}^{R}
      = \frac{2R^3}{3}.
\]
\textbf{Paso 3: Coordenada $\bar{y}$.}
\[
  \bar{y}=\frac{M_x}{A}
         = \frac{2R^3/3}{\pi R^2/2}
         = \boxed{\dfrac{4R}{3\pi}}.
\]
\end{myproof}
\end{example}

% -----------------------------------------------------------
\subsection{Momento de inercia}
% -----------------------------------------------------------

\begin{definition}[Momento de inercia de una lámina]
El \textbf{momento de inercia} de una lámina de densidad $\rho$ respecto
a un eje mide su resistencia a la rotación. Para los ejes coordenados:
\[
  I_x = \rho\int_a^b\frac{[f(x)]^3}{3}\,dx, \qquad
  I_y = \rho\int_a^b x^2 f(x)\,dx.
\]
El \textbf{momento polar} respecto al origen es $I_0=I_x+I_y$.
\end{definition}

\begin{rem}
El momento de inercia aparece en:
\begin{itemize}
  \item Energía cinética rotacional: $K=\frac{1}{2}I\omega^2$.
  \item Segunda ley de Newton para rotación: $\tau=I\alpha$.
\end{itemize}
A diferencia de la masa, el momento de inercia depende de \emph{cómo} está
distribuida la masa: puntos más lejanos al eje contribuyen más (factor
$r^2$). Por ello, reducir $I$ sin cambiar la masa (distribuyendo la masa
más cerca del eje) mejora la eficiencia en sistemas rotativos.
\end{rem}

\begin{example}[Momento de inercia de una varilla]
Una varilla delgada de longitud $L$ y masa $M$ gira alrededor de un
extremo. Calcule $I$.

\begin{myproof}
\textbf{Paso 1: Densidad lineal.}
$\lambda=M/L$. Un elemento $dx$ a distancia $x$ del eje tiene masa
$dm=\lambda\,dx$.

\textbf{Paso 2: Elemento diferencial.}
$dI = x^2\,dm = \lambda x^2\,dx$.

\textbf{Paso 3: Límites.} $x\in[0,L]$.

\textbf{Paso 4: Evaluación.}
\[
  I = \int_0^L\frac{M}{L} x^2\,dx
    = \frac{M}{L}\cdot\frac{L^3}{3}
    = \boxed{\dfrac{1}{3}ML^2}.
\]
\end{myproof}
\end{example}

\begin{example}[Momento de inercia respecto al centro]
Calcule el momento de inercia de la misma varilla respecto a su centro.

\begin{myproof}
\textbf{Paso 1: Sistema de coordenadas.}
Situamos el origen en el centro: $x\in[-L/2,L/2]$.

\textbf{Paso 2: Evaluación.}
\[
  I_{\text{centro}} = \int_{-L/2}^{L/2}\frac{M}{L}x^2\,dx
    = \frac{M}{L}\cdot\frac{2}{3}\left(\frac{L}{2}\right)^{\!3}
    = \frac{M}{L}\cdot\frac{L^3}{12}
    = \boxed{\dfrac{1}{12}ML^2}.
\]
Nótese que
$I_{\text{extremo}}=I_{\text{centro}}+M\!\left(\dfrac{L}{2}\right)^{\!2}
=\frac{1}{12}ML^2+\frac{1}{4}ML^2=\frac{1}{3}ML^2$,
lo que ilustra el \textbf{Teorema de Steiner (ejes paralelos):}
$I=I_{\text{cm}}+M d^2$, donde $d$ es la distancia entre los ejes.
\end{myproof}
\end{example}

% ============================================================
\section{Presión y fuerza hidrostática}
% ============================================================

\begin{definition}[Fuerza hidrostática sobre una placa plana]
La presión hidrostática a profundidad $h$ en un fluido de densidad $\rho$
es $P=\rho g h$. La fuerza total sobre una placa plana sumergida
verticalmente, con anchura $w(h)$ a profundidad $h$, entre profundidades
$h=a$ y $h=b$, es
\[
  F = \int_a^b\rho g h\,w(h)\,dh.
\]
\end{definition}

\begin{example}[Fuerza sobre una compuerta rectangular]
Una compuerta rectangular de $4\,\text{m}$ de ancho y $3\,\text{m}$ de
alto está sumergida verticalmente, con su borde superior a $2\,\text{m}$
de la superficie del agua ($\rho=1000\,\text{kg/m}^3$). Halle la fuerza
total sobre la compuerta.

\begin{myproof}
\textbf{Paso 1: Geometría y anchura.}
La compuerta ocupa profundidades $h\in[2,5]$ con anchura constante
$w(h)=4\,\text{m}$.

\textbf{Paso 2: Elemento diferencial.}
$dF = P(h)\cdot w(h)\,dh = (\rho g h)(4)\,dh = 39200\,h\,dh$.

\textbf{Paso 3: Límites.} $h\in[2,5]$.

\textbf{Paso 4: Evaluación.}
\[
  F = \int_2^5 39200\,h\,dh
    = 39200\left[\frac{h^2}{2}\right]_2^5
    = 39200\cdot\frac{25-4}{2}
    = 39200\cdot 10.5
    = \boxed{411600\,\text{N}\approx 411.6\,\text{kN}}.
\]
\end{myproof}
\end{example}

\section{Problemas propuestos}

% Requisitos mínimos en el preámbulo de tu documento:
% \usepackage{pgfplots}
% \pgfplotsset{compat=1.18}
% \usepackage{amsmath, amssymb}
% \newenvironment{prob}{\begin{trivlist}\item[\hskip \labelsep \textbf{Problema.}]}{\end{trivlist}} 
% (Ajusta la definición de `prob` según tu clase/base de plantilla)

Resuelva los siguientes problemas aplicando los protocolos establecidos en este capítulo. Revise cuidadosamente los ejemplos resueltos de cada tema antes de comenzar.

\begin{prob}
Calcule el área de la región delimitada por la gráfica de $f(x) = x^2$, el eje $x$ y la recta vertical $x = 2$.
\end{prob}

\begin{prob}
Determine el área bajo la curva $y = \sqrt{x}$ en el intervalo $[0, 4]$.
\end{prob}

\begin{prob}
Encuentre el área de la región delimitada por $y = x^3 - x^2 - 6x$ y el eje $x$.
\end{prob}

\begin{prob}
Encuentre el área de la región delimitada por las parábolas $y = x^2$ y $y = \sqrt{8x}$.
\end{prob}

\begin{prob}
Determine el área de la región acotada por las gráficas de $f(x) = x^2 + 2x - 3$ y $g(x) = -x + 1$.
\end{prob}

\begin{prob}
Calcule el área de la región encerrada entre $y = \sin x$ y $y = \frac{1}{2}$ para $0 \leq x \leq \frac{17\pi}{6}$.
\end{prob}

\begin{prob}
Calcule el área de la región entre la parábola $y^2 = 4x$ y la recta $4x - 3y = 4$.
\end{prob}

\begin{prob}
Encuentre el área de la región acotada por $x = 3 - y^2$ y $y = x - 1$.
\end{prob}

\begin{prob}
Determine el área de la región acotada por $x = y^2$ y $x = y + 2$.
\end{prob}

\begin{prob}
Considere la región $R$ limitada por las curvas $y=x+6$, $y=-x^2-6x-6$ y $x=0$. Plantee y evalúe la integral que permite calcular el área encerrada. Determine explícitamente los límites de integración y el diferencial correspondiente.
\begin{center}
\begin{tikzpicture}
  \begin{axis}[axis lines=middle, xlabel=$x$, ylabel=$y$, xmin=-4, xmax=1, ymin=-7, ymax=7, grid=major, width=8cm, height=7cm, axis on top]
    % Trazado de funciones con nombres para relleno
    \addplot[name path=f, domain=-3.5:0.5, samples=100, blue, thick] {x+6};
    \addplot[name path=g, domain=-3.5:0.5, samples=100, red, thick] {-x^2-6*x-6};
    % Relleno entre curvas solo en el dominio de interés [-3, 0]
    \addplot[purple!20, opacity=0.5] fill between[of=f and g, soft clip={domain=-3:0}];
    % Etiquetas
    \node at (axis cs:-2.5, 5.2) {\small $y=x+6$};
    \node at (axis cs:-2.2, -3.5) {\small $y=-x^2-6x-6$};
  \end{axis}
\end{tikzpicture}
\end{center}
\end{prob}

\begin{prob}
Calcule el volumen del sólido generado al girar la región bajo $y = \sqrt{x}$ desde $x = 0$ hasta $x = 4$ alrededor del eje $x$ (método de discos).
\end{prob}

\begin{prob}
Determine el volumen del sólido de revolución obtenido al girar la región acotada por $y = \cos x$ y el eje $x$ en $[0, \pi/2]$ alrededor del eje $x$.
\end{prob}

\begin{prob}
Calcule el volumen del sólido generado al girar la región acotada por $y = x^2$ y $y^2 = 8x$ en torno al eje $x$ (método de arandelas).
\end{prob}

\begin{prob}
Use el método de los discos (o arandelas) para calcular el volumen del sólido de revolución generado al girar el área encerrada por $y=2x+4$, $y=-2x^2-4x$ y $x=0$ respecto al eje $x$. Determine explícitamente los límites y el diferencial.
\begin{center}
\begin{tikzpicture}
  \begin{axis}[axis lines=middle, xlabel=$x$, ylabel=$y$, xmin=-2.5, xmax=0.5, ymin=-4, ymax=5, grid=major, width=8cm, height=7cm, axis on top]
    \addplot[name path=f, domain=-1.5:0.5, samples=100, blue, thick] {2*x+4};
    \addplot[name path=g, domain=-1.5:0.5, samples=100, red, thick] {-2*x^2-4*x};
    \addplot[purple!20, opacity=0.5] fill between[of=f and g, soft clip={domain=-1:0}];
    \node at (axis cs:-0.8, 4.2) {\small $y=2x+4$};
    \node at (axis cs:-0.8, 1.5) {\small $y=-2x^2-4x$};
  \end{axis}
\end{tikzpicture}
\end{center}
\end{prob}

\begin{prob}
La región acotada por $y = x^2$ y $y = 4x$ se gira alrededor del eje $x$. Use el método de las arandelas para calcular su volumen.
\end{prob}

\begin{prob}
Encuentre el volumen del sólido generado al girar la región semicircular $x = \sqrt{4-y^2}$ alrededor de la recta $x = -1$.
\end{prob}

\begin{prob}
Use el método de los cascarones cilíndricos para calcular el volumen del sólido generado al girar la región acotada por $y = 3 + 2x - x^2$ y $y = 0$ alrededor del eje $y$.
\end{prob}

\begin{prob}
Determine el volumen del sólido obtenido al girar la región acotada por $y = x$, $y = 0$ y $x = 2$ alrededor del eje $y$.
\end{prob}

\begin{prob}
Considere las funciones $f(x)=-x^2+4x-3$ y $g(x)=x^2-6x+9$.
\begin{enumerate} 
  \item Calcule el área entre las dos curvas.
  \item Calcule el volumen del sólido de revolución al girar la región respecto al eje $y=2$.
  \item Calcule el volumen del sólido de revolución al girar la región respecto al eje $x=3$.
\end{enumerate}
\end{prob}

\begin{prob}
Calcule el volumen del sólido generado al girar la región acotada por $y = e^{-x^2}$, $y = 0$, $x = 0$ y $x = 1$ alrededor del eje $y$.
\end{prob}

\begin{prob}
Calcule la longitud de la curva $y = x^{3/2}$ desde el punto $(1, 1)$ hasta el punto $(4, 8)$.
\end{prob}

\begin{prob}
Determine la longitud de arco de la curva $y = \ln(\cos x)$ desde $x = 0$ hasta $x = \pi/4$.
\end{prob}

\begin{prob}
Encuentre el área de la superficie generada al girar la curva $y = \sqrt{x}$ para $0 \leq x \leq 4$ alrededor del eje $x$.
\end{prob}

\begin{prob}
Demuestre que el área de la superficie de una esfera de radio $r$ es $4\pi r^2$ mediante el giro de una semicircunferencia.
\end{prob}

\begin{prob}
Un módulo espacial pesa 15 toneladas en la superficie terrestre. ¿Cuánto trabajo se requiere para propulsarlo a una altura de 800 millas? (Radio de la Tierra $\approx 4000$ millas).
\end{prob}

\begin{prob}
Un cable que pesa $2$ libras por pie se usa para levantar una carga de $200$ libras hasta la parte superior de un pozo que tiene $500$ pies de profundidad. ¿Cuánto trabajo se realiza?
\end{prob}

\begin{prob}
Un tanque tiene forma de cono circular recto invertido (altura $10$ pies, radio de la base $4$ pies). Si el tanque está lleno de agua, calcule el trabajo necesario para bombearla hasta el borde superior.
\end{prob}

\begin{prob}
Se requieren $0.09$ joules de trabajo para estirar un resorte desde $15$ a $17$ cm, y otros $0.10$ joules para estirarlo de $17$ a $20$ cm. Determine la constante del resorte y encuentre su longitud natural.
\end{prob}

\begin{prob}
Un tanque con forma de cono circular recto se usa para almacenar agua ($\rho=997\,\text{kg/m}^3$). El tanque tiene altura $10\,\text{m}$ y diámetro $16\,\text{m}$. Calcule el trabajo necesario para bombear el agua hasta una altura de $15\,\text{m}$ por encima del borde.
\end{prob}


\begin{prob}
Encuentre el centroide de la región acotada por las curvas $y = x^3$ y $y = \sqrt{x}$.
\end{prob}

\begin{prob}
Determine el centro de masa de una lámina con densidad $\rho(x,y) = xy$ acotada por $y = x^{2/3}$, $x = 8$ y el eje $x$.
\end{prob}

\begin{prob}
Encuentre el centroide del triángulo con vértices $(-a, 0)$, $(a, 0)$ y $(b, c)$.
\end{prob}

\begin{prob}
Calcule el momento de inercia $I_x$ de una lámina homogénea de densidad constante $k$ acotada por $y = x^2$ y $y = 4$.
\end{prob}

\begin{prob}
Determine el momento de inercia respecto al eje $y$ de un alambre delgado con forma de semicircunferencia de radio $a$.
\end{prob}


\begin{prob}
Una lámina de metal rectangular de $3 \times 4$ pies se sumerge verticalmente en agua a una profundidad de 6 pies (borde superior). Calcule la fuerza total ejercida por el fluido.
\end{prob}

\begin{prob}
Una compuerta en un canal tiene forma de trapezoide (3 pies de ancho abajo, 5 pies arriba, 2 pies de alto). Calcule la fuerza hidrostática si el nivel del agua coincide con la parte superior.
\end{prob}

\begin{prob}
Encuentre la fuerza de un fluido contra el extremo vertical de un depósito que tiene forma de región parabólica de 8 pies de ancho y 4 pies de profundidad.
\end{prob}

\begin{prob}[Desafío Putnam] Sea $f: [0, 1] \to \mathbb{R}$ una función continua. Se definen las siguientes integrales:
\[
I(f) = \int_{0}^{1} x^2 f(x) \, dx, \quad J(f) = \int_{0}^{1} x [f(x)]^2 \, dx
\]
Encuentre el valor máximo posible de la diferencia $I(f) - J(f)$ sobre todas las funciones $f$ posibles.
\end{prob}

\begin{prob}[Geometría de Curvas] Considere la hipocicloide de cuatro cúspides definida por la ecuación: $x^{2/3} + y^{2/3} = 1$.
\begin{enumerate} 
  \item Use la derivación implícita para encontrar $dy/dx$.
  \item Formule la integral de longitud de arco para el primer cuadrante. Observe que el integrando tiende a infinito cuando $x \to 0$.
  \item Demuestre que el perímetro total de la figura es exactamente $6$.
\end{enumerate}
\end{prob}



\begin{prob}[Integrabilidad y Composición] Encuentre dos funciones $f$ y $g$, ambas integrables en el intervalo $[0,1]$, tales que su composición $g \circ f$ \textbf{no} sea integrable en dicho intervalo. \textit{Sugerencia: Considere la función de Dirichlet y una modificación de la función de Thomae.}
\end{prob}

\begin{prob}[Aplicaciones de Densidad Variable] Una varilla tiene una densidad dada por $\delta(x) = e^{-x^2}$ para $x > 0$. 
\begin{enumerate} 
  \item Formule la integral que representa la masa total desde $x=0$ hasta un punto $x=a$.
  \item Dado que la antiderivada de $e^{-x^2}$ no es elemental, use el \textbf{método de bisección} o la \textbf{regla de Simpson} para determinar el punto de corte $a$ tal que la masa sea exactamente igual a $1$.
\end{enumerate}
\end{prob}




